\chapter{Operadores Derivados: NAND, NOR, XOR y XNOR}
Las ecuaciones obtenidas a partir de las Leyes de De Morgan demuestran que las operaciones básicas $\vee$ y $\wedge$ pueden ser expresadas íntegramente en términos de la negación de su operación dual. Esto motiva la definición de varios operadores lógicos fundamentales en sistemas digitales. Dos de ellos (NAND y NOR) son de gran relevancia por ser funcionalmente completos por sí solos, mientras que otros dos (XOR y XNOR) son esenciales para funciones aritméticas y de comprobación de paridad:

\begin{definicion}{Operador NAND (Barra de Sheffer)}{nand}
Denotado clásicamente con una flecha hacia arriba ($\uparrow$), se define como la negación del ínfimo.
\[ a \uparrow b \triangleq \neg (a \wedge b) = \neg a \vee \neg b \]
\end{definicion}

\begin{definicion}{Operador NOR (Flecha de Peirce)}{nor}
Denotado con una flecha hacia abajo ($\downarrow$), se define como la negación del supremo.
\[ a \downarrow b \triangleq \neg (a \vee b) = \neg a \wedge \neg b \]
\end{definicion}

\begin{definicion}{Operador XOR (O-exclusiva)}{xor}
Denotado con el símbolo de suma exclusiva ($\oplus$), evalúa a $\top$ cuando exactamente uno de los operandos es $\top$ y el otro $\bot$.
\[ a \oplus b \triangleq (a \wedge \neg b) \vee (\neg a \wedge b) \]
\end{definicion}

\begin{definicion}{Operador XNOR (No-O-exclusiva o Equivalencia)}{xnor}
Denotado frecuentemente con $\odot$ o $\leftrightarrow$, es la negación de la operación XOR y evalúa a $\top$ cuando ambos operandos son idénticos.
\[ a \odot b \triangleq \neg (a \oplus b) = (a \wedge b) \vee (\neg a \wedge \neg b) \]
\end{definicion}

\begin{definicion}{Generalización a $n$ variables de NAND y NOR}{gen_nand_nor}
A diferencia de los operadores $\vee$, $\wedge$ y $\oplus$, los operadores NAND ($\uparrow$) y NOR ($\downarrow$) \textbf{no son asociativos}. Sin embargo, debido a su inmensa importancia práctica en la construcción de circuitos digitales, se define convencionalmente su generalización a $n$ variables como la negación de la conjunción o disyunción múltiple, respectivamente:
\[ \text{NAND}(x_1, x_2, \dots, x_n) \triangleq \neg \left( \bigwedge_{i=1}^n x_i \right) \]
\[ \text{NOR}(x_1, x_2, \dots, x_n) \triangleq \neg \left( \bigvee_{i=1}^n x_i \right) \]
\end{definicion}

\section{Comportamiento de los Operadores Derivados}
Los operadores NAND ($\uparrow$) y NOR ($\downarrow$) presentan una serie de propiedades algebraicas particulares. Al ser operadores funcionalmente completos, permiten expresar cualquier otra operación booleana utilizando exclusivamente uno de ellos.

\begin{teorema}{Idempotencia cruzada (Generación de NOT)}{idemp_cruzada}
Operar un elemento consigo mismo usando NAND o NOR equivale a su complemento (negación):
\begin{align*}
    \forall a \in \mathbb{B}, \quad a \uparrow a &= \neg a \\
    \forall a \in \mathbb{B}, \quad a \downarrow a &= \neg a
\end{align*}
\end{teorema}
\desplegable{proof_idemp_cruzada}{
\begin{proof}[Demostración]
Para la operación NAND:
\begin{align*}
    a \uparrow a &\triangleq \neg (a \wedge a) & (\text{Definicion de NAND}) \\
                 &= \neg a & (Idemp_\wedge)
\end{align*}
Para la operación NOR:
\begin{align*}
    a \downarrow a &\triangleq \neg (a \vee a) & (\text{Definicion de NOR}) \\
                   &= \neg a & (Idemp_\vee)
\end{align*}
\end{proof}
}

\begin{teorema}{Generación del Ínfimo y Supremo (AND y OR)}{gen_inf_sup}
A partir de la propiedad anterior, podemos recuperar las operaciones básicas anidando las puertas consigo mismas:
\begin{align*}
    \forall a, b \in \mathbb{B}, \quad a \wedge b &= \neg (a \uparrow b) = (a \uparrow b) \uparrow (a \uparrow b) \\
    \forall a, b \in \mathbb{B}, \quad a \vee b &= \neg (a \downarrow b) = (a \downarrow b) \downarrow (a \downarrow b)
\end{align*}
\end{teorema}
\desplegable{proof_gen_inf_sup}{
\begin{proof}[Demostración]
Para la generación del ínfimo (AND):
\begin{align*}
    \neg(a \uparrow b) &\triangleq \neg(\neg(a \wedge b)) & (\text{Definicion de NAND}) \\
                       &= a \wedge b & (Involucion)
\end{align*}
Además, por la idempotencia cruzada demostrada anteriormente, $x \uparrow x = \neg x$, por tanto:
\[ \neg(a \uparrow b) = (a \uparrow b) \uparrow (a \uparrow b) \]
La demostración para la generación del supremo (OR) es idéntica por dualidad.
\end{proof}
}

\begin{teorema}{Generación cruzada (Leyes de De Morgan para NAND/NOR)}{gen_cruzada}
Podemos generar la operación opuesta (supremo desde NAND, e ínfimo desde NOR) negando previamente las entradas:
\begin{align*}
    \forall a, b \in \mathbb{B}, \quad a \vee b &= (\neg a) \uparrow (\neg b) = (a \uparrow a) \uparrow (b \uparrow b) \\
    \forall a, b \in \mathbb{B}, \quad a \wedge b &= (\neg a) \downarrow (\neg b) = (a \downarrow a) \downarrow (b \downarrow b)
\end{align*}
\end{teorema}
\desplegable{proof_gen_cruzada}{
\begin{proof}[Demostración]
Para el supremo (OR):
\begin{align*}
    (\neg a) \uparrow (\neg b) &\triangleq \neg (\neg a \wedge \neg b) & (\text{Definicion de NAND}) \\
                               &= \neg (\neg (a \vee b)) & (Mor_\vee) \\
                               &= a \vee b & (Involucion)
\end{align*}
La demostración para el ínfimo (AND) sigue los mismos pasos de manera dual, aplicando $Mor_\wedge$.
\end{proof}
}

\begin{teorema}{Conmutatividad}{conmut_deriv}
Al igual que sus operaciones base, ambos operadores son perfectamente conmutativos:
\begin{align*}
    \forall a, b \in \mathbb{B}, \quad a \uparrow b &= b \uparrow a \\
    \forall a, b \in \mathbb{B}, \quad a \downarrow b &= b \downarrow a
\end{align*}
\end{teorema}
\desplegable{proof_conmut_deriv}{
\begin{proof}[Demostración]
Para la operación NAND:
\begin{align*}
    a \uparrow b &\triangleq \neg (a \wedge b) & (\text{Definicion de NAND}) \\
                 &= \neg (b \wedge a) & (Comm_\wedge) \\
                 &\triangleq b \uparrow a & (\text{Definicion de NAND})
\end{align*}
Para la operación NOR, es análogo aplicando $Comm_\vee$.
\end{proof}
}

\begin{teorema}{Inexistencia de Elemento Neutro}{no_neutro}
No existe ningún elemento neutro para las operaciones NAND ni NOR en un álgebra de Boole general.
\end{teorema}
\desplegable{proof_no_neutro}{
\begin{proof}[Demostración]
Si existiera un neutro $e$ para la operación NAND, debería cumplirse que $\forall a, a \uparrow e = a$, es decir, $\neg(a \wedge e) = a$. Si probamos con $e=\top$, obtenemos $\neg a = a$, lo cual obliga al colapso en un álgebra trivial. Si probamos con $e=\bot$, obtenemos $\neg \bot = a \implies \top = a$, lo cual obviamente no se cumple para cualquier elemento $a$. Lo mismo aplica a la operación NOR.
\end{proof}
}

\begin{teorema}{Comportamiento con las constantes (Fijación y Absorción)}{constantes_deriv}
Fijar una constante específica en uno de los operandos genera directamente la negación, mientras que usar la constante opuesta actúa como un pseudo-elemento absorbente (devolviendo un valor constante inalterable por $a$):
\begin{align*}
    \text{Inversión: } & a \uparrow \top = \neg a \qquad & a \downarrow \bot &= \neg a \\
    \text{Absorción: } & a \uparrow \bot = \top \qquad & a \downarrow \top &= \bot
\end{align*}
\end{teorema}
\desplegable{proof_constantes_deriv}{
\begin{proof}[Demostración]
Para la inversión:
\begin{align*}
    a \uparrow \top &= \neg(a \wedge \top) = \neg a & (ElemNeu_\wedge) \\
    a \downarrow \bot &= \neg(a \vee \bot) = \neg a & (ElemNeu_\vee)
\end{align*}
Para la absorción:
\begin{align*}
    a \uparrow \bot &= \neg(a \wedge \bot) = \neg \bot = \top & (Abs_\bot) \\
    a \downarrow \top &= \neg(a \vee \top) = \neg \top = \bot & (Abs_\top)
\end{align*}
\end{proof}
}

\begin{teorema}{Ausencia de Asociatividad}{no_asoc_deriv}
A diferencia del supremo ($\vee$) y el ínfimo ($\wedge$), las operaciones NAND y NOR son positivamente NO asociativas:
\begin{align*}
    (a \uparrow b) \uparrow c &\neq a \uparrow (b \uparrow c) \\
    (a \downarrow b) \downarrow c &\neq a \downarrow (b \downarrow c)
\end{align*}
\end{teorema}
\desplegable{proof_no_asoc_deriv}{
\begin{proof}[Demostración]
Desarrollando el lado izquierdo para NAND:
\[ (a \uparrow b) \uparrow c = \neg ( (\neg (a \wedge b)) \wedge c ) = (a \wedge b) \vee \neg c \]
Desarrollando el lado derecho:
\[ a \uparrow (b \uparrow c) = \neg ( a \wedge (\neg (b \wedge c)) ) = \neg a \vee (b \wedge c) \]
Resulta evidente que $(a \wedge b) \vee \neg c \neq \neg a \vee (b \wedge c)$ para combinaciones arbitrarias de variables. El mismo razonamiento aplica de manera estricta y análoga para demostrar la carencia de asociatividad en la operación NOR ($\downarrow$).
\end{proof}
}

\begin{definicion}{NAND y NOR de $n$ entradas}{def_n_entradas}
Dado que las puertas NAND y NOR físicas a menudo tienen más de dos entradas, se definen algebraicamente para múltiples entradas como la negación de la conjunción o disyunción de todas ellas:
\begin{align*}
    \uparrow(x_1, x_2, \dots, x_n) &\triangleq \neg \left( \bigwedge_{i=1}^n x_i \right) \\
    \downarrow(x_1, x_2, \dots, x_n) &\triangleq \neg \left( \bigvee_{i=1}^n x_i \right)
\end{align*}
En particular, para el caso de 3 entradas que estudiaremos a continuación:
\begin{align*}
    \uparrow(a,b,c) &\triangleq \neg(a \wedge b \wedge c) \\
    \downarrow(a,b,c) &\triangleq \neg(a \vee b \vee c)
\end{align*}
\end{definicion}

\begin{teorema}{NAND/NOR múltiple vs agrupación binaria}{multiple_vs_binaria}
Como consecuencia directa de su falta de asociatividad, una operación NAND o NOR de 3 entradas no es equivalente a la agrupación secuencial en cascada de operaciones de 2 entradas:
\begin{align*}
    \uparrow(a,b,c) &\neq (a \uparrow b) \uparrow c \qquad \text{y} \qquad \uparrow(a,b,c) \neq a \uparrow (b \uparrow c) \\
    \downarrow(a,b,c) &\neq (a \downarrow b) \downarrow c \qquad \text{y} \qquad \downarrow(a,b,c) \neq a \downarrow (b \downarrow c)
\end{align*}
\end{teorema}
\desplegable{proof_multiple_vs_binaria}{
\begin{proof}[Demostración]
Para la NAND de 3 entradas tenemos, por definición y leyes de De Morgan:
\[ \uparrow(a,b,c) \triangleq \neg(a \wedge b \wedge c) = \neg a \vee \neg b \vee \neg c \]
Sin embargo, la agrupación de dos en dos evaluada anteriormente daba:
\[ (a \uparrow b) \uparrow c = (a \wedge b) \vee \neg c \]
Evidentemente, $\neg a \vee \neg b \vee \neg c \neq (a \wedge b) \vee \neg c$. Lo mismo aplica a las agrupaciones derechas y a las operaciones NOR equivalentes.
\end{proof}
}

\begin{teorema}{Extensión del Principio de Dualidad (NAND y NOR)}{dualidad_nand_nor}
La inclusión de los operadores derivados expande el Principio de Dualidad establecido en los postulados iniciales. La expresión dual de cualquier teorema o identidad que contenga operaciones NAND o NOR se obtiene intercambiando los operadores $\uparrow$ y $\downarrow$ (además de los ya conocidos $\vee \leftrightarrow \wedge$ y $\bot \leftrightarrow \top$).
\end{teorema}
\subsection{Comportamiento de los Operadores XOR y XNOR}
A diferencia de los operadores NAND y NOR, que destacan por su universalidad funcional pero carecen de propiedades algebraicas deseables (como asociatividad o elemento neutro), los operadores XOR ($\oplus$) y XNOR ($\odot$) exhiben una rica estructura algebraica. A continuación demostraremos estas propiedades.

\begin{teorema}{Conmutatividad}{conmut_xor}
Ambos operadores son conmutativos:
\[ a \oplus b = b \oplus a \qquad \text{y} \qquad a \odot b = b \odot a \]
\end{teorema}
\desplegable{proof_conmut_xor}{
\begin{proof}[Demostración]
Para XOR:
\begin{align*}
    a \oplus b &\triangleq (a \wedge \neg b) \vee (\neg a \wedge b) & (\text{Definición de XOR}) \\
               &= (\neg a \wedge b) \vee (a \wedge \neg b) & (Comm_\vee) \\
               &= (b \wedge \neg a) \vee (\neg b \wedge a) & (Comm_\wedge) \\
               &\triangleq b \oplus a & (\text{Definición de XOR})
\end{align*}
La demostración para XNOR sigue pasos idénticos aprovechando la conmutatividad del ínfimo y el supremo.
\end{proof}
}

\begin{teorema}{Elementos Neutros e Inversores}{neutro_xor}
El elemento $\bot$ actúa como neutro para la XOR, y $\top$ actúa como inversor. De manera dual, $\top$ es el neutro de la XNOR, y $\bot$ actúa como inversor:
\begin{align*}
    a \oplus \bot &= a & a \oplus \top &= \neg a \\
    a \odot \top &= a & a \odot \bot &= \neg a
\end{align*}
\end{teorema}
\desplegable{proof_neutro_xor}{
\begin{proof}[Demostración para XOR]
Para el elemento $\bot$ (neutro):
\begin{align*}
    a \oplus \bot &\triangleq (a \wedge \neg \bot) \vee (\neg a \wedge \bot) & (\text{Definición}) \\
                  &= (a \wedge \top) \vee \bot & (Comp_{inv} \text{ y } Fij_\wedge) \\
                  &= a \vee \bot & (ElemNeu_\wedge) \\
                  &= a & (ElemNeu_\vee)
\end{align*}
Para el elemento $\top$ (inversor):
\begin{align*}
    a \oplus \top &\triangleq (a \wedge \neg \top) \vee (\neg a \wedge \top) & (\text{Definición}) \\
                  &= (a \wedge \bot) \vee \neg a & (Comp_{inv} \text{ y } ElemNeu_\wedge) \\
                  &= \bot \vee \neg a & (Fij_\wedge) \\
                  &= \neg a & (ElemNeu_\vee)
\end{align*}
Las pruebas para XNOR son totalmente duales.
\end{proof}
}

\begin{teorema}{Elemento Inverso de sí mismo (Grupo Abeliano)}{idemp_nula_xor}
La combinación de un elemento consigo mismo produce una anulación (devuelve el neutro de la operación correspondiente), actuando cada elemento como su propio inverso:
\[ a \oplus a = \bot \qquad \text{y} \qquad a \odot a = \top \]
\end{teorema}
\desplegable{proof_idemp_nula_xor}{
\begin{proof}[Demostración]
Para XOR:
\begin{align*}
    a \oplus a &\triangleq (a \wedge \neg a) \vee (\neg a \wedge a) \\
               &= \bot \vee \bot & (Comp_\wedge \text{ y } Comm_\wedge) \\
               &= \bot & (Idemp_\vee)
\end{align*}
Para XNOR:
\begin{align*}
    a \odot a &\triangleq (a \wedge a) \vee (\neg a \wedge \neg a) \\
              &= a \vee \neg a & (Idemp_\wedge) \\
              &= \top & (Comp_\vee)
\end{align*}
\end{proof}
}

\begin{teorema}{Propiedades de Negación}{negacion_xor}
Negar cualquiera de las entradas de forma independiente equivale a negar la operación completa, lo que a su vez alterna entre XOR y XNOR:
\begin{align*}
    \neg (a \oplus b) &= \neg a \oplus b = a \oplus \neg b = a \odot b \\
    \neg (a \odot b) &= \neg a \odot b = a \odot \neg b = a \oplus b
\end{align*}
\end{teorema}
\desplegable{proof_negacion_xor}{
\begin{proof}[Demostración de $a \oplus \neg b = \neg(a \oplus b)$]
\begin{align*}
    a \oplus \neg b &\triangleq (a \wedge \neg(\neg b)) \vee (\neg a \wedge \neg b) \\
                    &= (a \wedge b) \vee (\neg a \wedge \neg b) & (Involucion) \\
                    &\triangleq a \odot b & (\text{Definición de XNOR})
\end{align*}
Sabiendo por definición que $a \odot b \triangleq \neg(a \oplus b)$, se concluye de forma inmediata que $a \oplus \neg b = \neg (a \oplus b)$.
\end{proof}
}

\begin{teorema}{Asociatividad}{asoc_xor}
Tanto XOR como XNOR son operadores algebraicamente asociativos:
\[ (a \oplus b) \oplus c = a \oplus (b \oplus c) \qquad \text{y} \qquad (a \odot b) \odot c = a \odot (b \odot c) \]
\end{teorema}
\desplegable{proof_asoc_xor}{
\begin{proof}[Demostración para XOR]
Primero evaluaremos el miembro izquierdo $(a \oplus b) \oplus c$. Llamaremos $X = a \oplus b$.
\begin{align*}
    X \oplus c &\triangleq (X \wedge \neg c) \vee (\neg X \wedge c) \\
               &= \big( ((a \wedge \neg b) \vee (\neg a \wedge b)) \wedge \neg c \big) \vee \big( \neg ((a \wedge \neg b) \vee (\neg a \wedge b)) \wedge c \big) \\
               &= \big( ((a \wedge \neg b) \vee (\neg a \wedge b)) \wedge \neg c \big) \vee \big( (a \odot b) \wedge c \big) \quad (\text{Definición de XNOR}) \\
               &= \big( (a \wedge \neg b \wedge \neg c) \vee (\neg a \wedge b \wedge \neg c) \big) \vee \big( ((a \wedge b) \vee (\neg a \wedge \neg b)) \wedge c \big) \quad (Dist_\wedge) \\
               &= (a \wedge \neg b \wedge \neg c) \vee (\neg a \wedge b \wedge \neg c) \\
               &\quad \vee (a \wedge b \wedge c) \vee (\neg a \wedge \neg b \wedge c) \quad (Dist_\wedge)
\end{align*}
Ahora evaluaremos el miembro derecho $a \oplus (b \oplus c)$. Llamaremos $Y = b \oplus c$.
\begin{align*}
    a \oplus Y &\triangleq (a \wedge \neg Y) \vee (\neg a \wedge Y) \\
               &= (a \wedge (b \odot c)) \vee (\neg a \wedge ((b \wedge \neg c) \vee (\neg b \wedge c))) \\
               &= (a \wedge ((b \wedge c) \vee (\neg b \wedge \neg c))) \vee (\neg a \wedge b \wedge \neg c) \vee (\neg a \wedge \neg b \wedge c) \\
               &= (a \wedge b \wedge c) \vee (a \wedge \neg b \wedge \neg c) \\
               &\quad \vee (\neg a \wedge b \wedge \neg c) \vee (\neg a \wedge \neg b \wedge c)
\end{align*}
Como podemos observar, ambas expansiones resultan exactamente en los mismos cuatro minitérminos. Reordenándolos por conmutatividad ($Comm_\vee$) demostramos que son idénticos.
\end{proof}
}

\begin{teorema}{Generalización n-aria (XOR y XNOR)}{gen_xor}
Dado que ambos operadores han demostrado ser asociativos y conmutativos, es posible omitir los paréntesis y generalizar la operación a un número arbitrario $n$ de variables. Se denotan mediante los operadores de sumatoria y productorio modificados:
\[ \bigoplus_{i=1}^{n} x_i = x_1 \oplus x_2 \oplus \dots \oplus x_n \]
\[ \bigodot_{i=1}^{n} x_i = x_1 \odot x_2 \odot \dots \odot x_n \]

\textbf{Propiedad de Paridad de la XOR (Suma Módulo 2):}
La operación $\bigoplus_{i=1}^{n} x_i$ es conocida matemáticamente como la \textit{función de paridad impar}. Su valor será $\top$ (o 1) si y solo si un número impar de las variables de entrada $x_i$ tienen valor $\top$.

\textbf{Relación entre XOR y XNOR para $n$ variables:}
Dado que $x \odot y = x \oplus y \oplus \top$, cada aplicación sucesiva del operador $\odot$ es matemáticamente equivalente a aplicar una suma $\oplus$ y concatenar una constante $\top$. Puesto que la expresión $\bigodot_{i=1}^{n} x_i$ entrelaza a sus operandos mediante exactamente $n-1$ operadores de equivalencia, se deduce rigurosamente que:
\[ \bigodot_{i=1}^{n} x_i = \left( \bigoplus_{i=1}^{n} x_i \right) \oplus \underbrace{\top \oplus \top \dots \oplus \top}_{n-1 \text{ veces}} \]
Por la propia idempotencia nula de la suma booleana exclusiva ($x \oplus x = \bot$), la suma de $n-1$ constantes $\top$ será igual a $\bot$ si $n-1$ es par (lo cual ocurre cuando $n$ es impar), y será igual a $\top$ si $n-1$ es impar (cuando $n$ es par). Por lo tanto:
\[ 
\bigodot_{i=1}^{n} x_i = \begin{cases} 
\bigoplus_{i=1}^{n} x_i & \text{si } n \text{ es impar (es la misma función)} \\
\neg \left( \bigoplus_{i=1}^{n} x_i \right) & \text{si } n \text{ es par (son funciones negadas)}
\end{cases} 
\]
\end{teorema}

\begin{teorema}{Distributividad (AND sobre XOR y OR sobre XNOR)}{dist_xor}
El producto (AND) se distribuye sobre la suma exclusiva (XOR), y dualmente, la suma (OR) se distribuye sobre la equivalencia (XNOR):
\[ a \wedge (b \oplus c) = (a \wedge b) \oplus (a \wedge c) \]
\[ a \vee (b \odot c) = (a \vee b) \odot (a \vee c) \]
\end{teorema}
\desplegable{proof_dist_xor}{
\begin{proof}[Demostración de AND sobre XOR]
Desarrollando el lado derecho (RHS):
\begin{align*}
    (a \wedge b) \oplus (a \wedge c) &\triangleq ((a \wedge b) \wedge \neg(a \wedge c)) \vee (\neg(a \wedge b) \wedge (a \wedge c)) \\
                                     &= (a \wedge b \wedge (\neg a \vee \neg c)) \vee ((\neg a \vee \neg b) \wedge a \wedge c) \quad (Mor_\wedge) \\
                                     &= ((a \wedge b \wedge \neg a) \vee (a \wedge b \wedge \neg c)) \\
                                     &\quad \vee ((a \wedge c \wedge \neg a) \vee (a \wedge c \wedge \neg b)) \quad (Dist_\wedge) \\
                                     &= (\bot \vee (a \wedge b \wedge \neg c)) \vee (\bot \vee (a \wedge \neg b \wedge c)) \quad (Comp_\wedge \text{ y } Fij_\wedge) \\
                                     &= (a \wedge b \wedge \neg c) \vee (a \wedge \neg b \wedge c) \quad (ElemNeu_\vee) \\
                                     &= a \wedge ((b \wedge \neg c) \vee (\neg b \wedge c)) \quad (Dist_\wedge \text{ a la inversa}) \\
                                     &\triangleq a \wedge (b \oplus c) \quad (\text{Def. XOR})
\end{align*}
Lo cual demuestra la igualdad. La demostración de la distributividad para XNOR es rigurosamente dual.
\end{proof}
}

\begin{teorema}{Extensión del Principio de Dualidad (XOR y XNOR)}{dualidad_xor_xnor}
De forma análoga a la relación entre NAND y NOR, los operadores XOR y XNOR son mutuamente duales. Para obtener la expresión dual de cualquier proposición que involucre estos operadores, se deben intercambiar $\oplus$ y $\odot$, manteniendo las reglas de dualidad estándar para los demás elementos y constantes.
\end{teorema}

\section{Estructuras Algebraicas Superiores}

Las propiedades demostradas anteriormente para los operadores XOR ($\oplus$) y AND ($\wedge$) permiten dotar al conjunto $B$ de estructuras algebraicas más ricas y estándar dentro del álgebra abstracta.

\begin{teorema}{Estructura de Grupo Conmutativo}{grupo_xor}
El par $(B, \oplus)$ forma un \textbf{grupo abeliano} (conmutativo), satisfaciendo:
\begin{enumerate}
    \item \textbf{Clausura:} $\forall a,b \in B, a \oplus b \in B$.
    \item \textbf{Asociatividad:} $a \oplus (b \oplus c) = (a \oplus b) \oplus c$.
    \item \textbf{Elemento neutro:} Existe $\bot \in B$ tal que $a \oplus \bot = a$.
    \item \textbf{Elemento inverso (Idempotencia aditiva):} Todo elemento es su propio inverso, ya que $a \oplus a = \bot$.
    \item \textbf{Conmutatividad:} $a \oplus b = b \oplus a$.
\end{enumerate}
\end{teorema}

\begin{teorema}{Estructura de Anillo Conmutativo (Anillo Booleano)}{anillo_booleano}
La terna $(B, \oplus, \wedge)$ forma un \textbf{anillo conmutativo con elemento unidad} (comúnmente denominado Anillo Booleano), donde el XOR actúa como la suma del anillo y el AND como el producto. Se satisfacen todas las propiedades requeridas:
\begin{enumerate}
    \item $(B, \oplus)$ es un grupo abeliano (demostrado arriba).
    \item $(B, \wedge)$ es un monoide conmutativo (asociativo, conmutativo, y con elemento neutro $\top$).
    \item \textbf{Distributividad:} El producto ($\wedge$) se distribuye sobre la suma ($\oplus$), es decir, $a \wedge (b \oplus c) = (a \wedge b) \oplus (a \wedge c)$.
\end{enumerate}
\end{teorema}

